\subsection{Recent paper-wise analysis: 2023--2026}

\paragraph{Liu et al. (2023).}
Liu et al. sought to improve electricity-theft detection by learning differences among customer representations rather than relying only on independently encoded histories \cite{Liu2023ContrastiveETD}. On SGCC customer histories with customer-level normal/theft targets, they combined a supervised contrastive encoder, largest-triangle-three-buckets augmentation, and a classification objective. The paper evaluates AUC and recall and reports favourable results under its SGCC protocol. Its relevance is the explicit use of learned customer representations. The verified evidence does not establish a matched direct-distribution control, and its customer partition is not the governed AMP split; its scores are therefore contextual rather than directly comparable.

\paragraph{Pamir et al. (2023).}
Pamir et al. addressed dimensionality and class imbalance in SGCC electricity-theft classification \cite{Pamir2023EnergyOptimization}. Their pipeline combines cost-sensitive learning, a gated convolutional autoencoder, an LSTM classifier, and salp-swarm hyperparameter tuning. The article states 42,372 customer histories, 3,615 theft customers, and an 80/20 train/test split. It reports 99.45\% precision, 95.93\% F1, 92.25\% accuracy, and 71.13\% AUC. The unusual divergence among these metrics reinforces the need to preserve source definitions and prevalence. Extensive tuning and the absence of AMP's frozen cohorts prevent protocol-level comparison.

\paragraph{Khan et al. (2023).}
Khan et al. investigated an AlexNet-derived deep classifier for customer-level electricity-theft detection on CNGC/SGCC consumption histories \cite{Khan2023AlexNet}. The verified package records imbalance treatment and evaluation with accuracy, precision, recall, F1, and AUC. The study supports the conclusion that deep full-history models can extract discriminative consumption patterns. Exact protocol equivalence with AMP was not established: its sampling, model tuning, and split differ, and no matched direct-versus-latent distribution comparison is reported. Accordingly, the paper informs the temporal-classifier context rather than a numerical ranking.

\paragraph{Huang et al. (2024).}
Huang et al. proposed dual-time feature fusion to combine complementary temporal views of SGCC daily customer histories \cite{Huang2024DualTime}. Under the paper's customer-supervised split, the reported values are 0.947 accuracy, 0.932 precision, 0.964 recall, 0.948 F1, and 0.986 AUC. The findings show that multi-scale temporal organization can be useful for theft-associated customer classification. Average Precision was not reported in the verified comparison record, and the preprocessing and split are not AMP-governed. The model also studies temporal fusion rather than the specific question of preserving the distribution of unlabeled windows.

\paragraph{Zhuang et al. (2024).}
Zhuang et al. developed a dynamic generative residual graph convolutional network for electricity-theft detection on SGCC customer histories \cite{Zhuang2024GraphETD}. The objective is to exploit relational or graph-structured dependencies while augmenting the representation through a generative component. The article reports classification metrics under its own preprocessing and partition. The approach broadens the temporal deep-learning family beyond sequence-only models, but a precise point finding suitable for transcription was not established by the Stage 29 package. For AMP's question, graph construction and generative augmentation change several factors simultaneously; the verified evidence does not isolate aggregation or provide a matched direct-feature control.

\paragraph{Shahzadi et al. (2024).}
Shahzadi et al. proposed a data-driven artificial-intelligence pipeline for energy-theft detection in urbanized smart grids using SGCC customer histories \cite{Shahzadi2024Urbanized}. The verified record identifies a deep-learning approach with imbalance handling and reports the standard accuracy, precision, recall, F1, and AUC family under a paper-specific partition. The study is relevant as a recent supervised SGCC benchmark. The verified package establishes the metric family but not an exact point finding suitable for transcription, so no value is reproduced here. Optimization, preprocessing, and evaluation differences limit direct comparison with AMP-LDP.

\paragraph{Tursunboev et al. (2024).}
Tursunboev et al. combined multi-objective evolutionary optimization with hybrid deep learning for energy-theft detection on SGCC customer histories \cite{Tursunboev2024MultiObjective}. The objective is to improve supervised customer classification by optimizing competing model goals. The article reports customer-level classification metrics within its internally optimized protocol. It illustrates the performance-oriented search direction in recent electricity-theft research, while the Stage 29 package does not establish an exact point finding for transcription. The search budget, optimization objectives, and lack of a separately frozen external cohort differ from AMP's governance; the verified package therefore supports only a structural comparison, not a transfer of headline scores.

\paragraph{Wang et al. (2024).}
Wang et al. studied a two-stage approach intended to generalize electricity-theft detection to new regions \cite{Wang2024NewRegions}. The environment contains source and new-region theft data, and the evaluation focuses on performance under regional transfer. The verified finding is protocol-level evidence from a cross-region assessment; an exact point result was not established for this package. This is important because it targets a limitation that AMP does not resolve: evidence from one source does not establish utility- or region-level generalization. The method and metrics are tied to the paper's cross-region protocol, which differs from AMP's within-SGCC governed cohorts.

\paragraph{Iftikhar et al. (2024).}
Iftikhar et al. addressed high dimensionality and imbalance in CNGC/SGCC customer histories using K-means SMOTE with an MLP--GRU classifier \cite{Iftikhar2024MLPGRU}. They report precision, recall, F1, accuracy, MCC, ROC-AUC, and PR-AUC across 50/50, 75/25, and 90/10 train/test ratios; the Stage 29 record notes a displayed best AUC of 0.95. The multi-ratio analysis demonstrates split sensitivity. Synthetic balancing and the absence of fixed cohort identities make the protocol unlike AMP's natural-prevalence, governed evaluation, so the value is not used as a direct comparator.

\paragraph{Saqib et al. (2024).}
Saqib et al. examined deep-learning-based electricity-theft prediction in a non-smart-grid billing or consumption environment with a customer theft target \cite{Saqib2024NonSmartGrid}. The article evaluates a supervised deep predictor with classification metrics under its own split. Its principal contribution to this review is environmental breadth: theft-detection systems may operate with data and infrastructure different from SGCC smart-meter histories. A precise point finding was not established by Stage 29, and those differences in acquisition setting and feature availability also limit comparability; no numerical result is inferred.

\paragraph{Elshennawy et al. (2025).}
Elshennawy et al. proposed a hybrid deep-learning classifier for electricity-theft detection on SGCC customer histories \cite{Elshennawy2025EfficientETD}. The study uses customer supervision and reports accuracy up to 94.7\%, together with other classification measures, under its article-specific protocol. It provides recent evidence that hybrid deep representations can discriminate SGCC customer labels. The verified record does not establish alignment with AMP's preprocessing, prevalence, split membership, or paired controls. The reported accuracy is therefore retained as a source-specific result rather than evidence of superiority over or inferiority to AMP-LDP.

\paragraph{Iqbal et al. (2025).}
Iqbal et al. investigated transfer learning for temporal electricity sequences using CNN--LSTM and BiGRU--BiLSTM models \cite{Iqbal2025TransferTemporal}. The customer-level evaluation reports precision, recall, F1, accuracy, and ROC-AUC, with ROC-AUC up to 0.89 in the article protocol. The work shows how temporal representations can be transferred rather than trained only within one target sample. The construction of source and target data differs from AMP's fixed DEV\_TRAIN, DEV\_VAL, CONFIRMATION, and final roles; the result therefore provides methodological context but not a protocol-equivalent benchmark.

\paragraph{Yang et al. (2026).}
Yang et al. combined a temporal convolutional network with Bayesian hyperparameter optimization and post-hoc explainability for electricity-theft classification \cite{Liupeng2026ExplainableETD}. The paper uses customer consumption time series and reports 98.04\% accuracy under its own stratified protocol. Its objective spans temporal feature learning, parameter search, and explanation through feature-attribution methods. Accuracy is sensitive to prevalence, and the cohort and tuning procedure are not AMP-equivalent. The paper is relevant to temporal convolution and interpretability, but its result cannot be inserted into an AMP leaderboard.

\paragraph{Yang et al. (2024).}
Zhengqiang Yang et al. proposed a self-decision ant-colony clustering algorithm for electricity-theft detection and assessed it on SGCC and Smart* household data \cite{Yang2024AntColony}. The method is an unsupervised clustering approach followed by evaluation against customer or household labels, offering a non-neural alternative to supervised deep models. The verified finding is that the method received dataset-specific classification assessment; the package does not establish a single point estimate suitable for AMP comparison. Use of two datasets and a clustering objective changes the scientific question relative to calibrated customer-probability modelling.

\paragraph{Gunduz and Das (2024).}
Gunduz and Das investigated a hybrid CNN-based security approach for detecting energy theft from consumption patterns \cite{Gunduz2024HybridCNN}. The verified record identifies constructed or injected attack patterns, generative data treatment, and evaluation through accuracy, precision, recall, F1, and ROC analysis. The study demonstrates a controlled attack-pattern setting in which convolutional models can be tested, although the verified package does not establish an exact point finding for transcription. Its labels are not equivalent to SGCC's inspected customer flag: synthetic or constructed attacks provide temporal and behavioural assumptions absent from the AMP data. This protocol difference is the central limitation for comparison.
